\documentclass[12pt]{article}
\usepackage[UTF8]{inputenc}
\usepackage{thaisupp}
\usepackage{geometry}
\geometry{a4paper, margin=2.5cm}
\begin{document}
\begin{center}
{\LARGE\textbf{ฟิสิกส์ ม.ปลาย — Part 1}}\\[0.3em]
{\Large\textbf{Lesson 9: คลื่น (Extra) (Extra)}}\\[0.5em]
\rule{0.8\linewidth}{0.5pt}
\end{center}
\section*{คำถาม In-Class}
\begin{enumerate}
\item \textbf{ข้อใดไม่ใช่คลื่นกล}
\begin{enumerate}
\item ก) คลื่นเสียง
\item ข) คลื่นแม่เหล็กไฟฟ้า
\item ค) คลื่นผิวน้ำ
\item ง) คลื่นในสปริง
\end{enumerate}
\item \textbf{v = fλ คือสมการอะไร}
\begin{enumerate}
\item ก) กฎการสะท้อน
\item ข) สมการความสัมพันธ์ของคลื่น
\item ค) กฎสเนลล์
\item ง) สมการพลังงาน
\end{enumerate}
\item \textbf{อัตราเร็วเสียงในอากาศที่อุณหภูมิ 25°C ประมาณเท่าใด}
\begin{enumerate}
\item ก) 300 m/s
\item ข) 340 m/s
\item ค) 380 m/s
\item ง) 400 m/s
\end{enumerate}
\item \textbf{เมื่อคลื่นเข้าหาฝั่งแล้วสะท้อนกลับ ข้อใดถูกต้องเกี่ยวกับความถี่}
\begin{enumerate}
\item ก) ความถี่เพิ่มขึ้น
\item ข) ความถี่ลดลง
\item ค) ความถี่คงที่
\item ง) ขึ้นอยู่กับมุมสะท้อน
\end{enumerate}
\item \textbf{การหักเหของคลื่นเกิดขึ้นเมื่อคลื่นเคลื่อนที่จากตัวกลางหนึ่งไปยังอีกตัวกลางหนึ่งที่มีสมบัติอย่างไร}
\begin{enumerate}
\item ก) มีความหนาแน่นเท่ากัน
\item ข) มีความหนาแน่นต่างกัน
\item ค) มีอุณหภูมิเท่ากัน
\item ง) มีความดันเท่ากัน
\end{enumerate}
\item \textbf{คลื่นน้ำมีความยาวคลื่น 2 m อัตราเร็ว 4 m/s ความถี่และคาบเวลาเป็นเท่าใด}
\begin{enumerate}
\item ก) f = 2 Hz, T = 0.5 s
\item ข) f = 2 Hz, T = 0.5 s
\item ค) f = 8 Hz, T = 0.125 s
\item ง) f = 6 Hz, T = 0.167 s
\end{enumerate}
\item \textbf{ข้อใดต่อไปนี้เป็นตัวอย่างของคลื่นตามยาว}
\begin{enumerate}
\item ก) คลื่นแม่เหล็กไฟฟ้า
\item ข) คลื่นผิวน้ำ
\item ค) คลื่นเสียง
\item ง) คลื่นในเส้นเชือก
\end{enumerate}
\item \textbf{ในการทดลองแทรกสอดของคลื่นแสงช่องคู่ ระยะระหว่างปฏิบัพที่ 1 และปฏิบัพที่ 5 บนฉายวัดได้ 12 mm ความยาวคลื่นของแสงเป็นเท่าใด (d = 0.2 mm, D = 1 m)}
\begin{enumerate}
\item ก) 400 nm
\item ข) 500 nm
\item ค) 600 nm
\item ง) 700 nm
\end{enumerate}
\item \textbf{คลื่นกลต้องมีสมบัติอย่างไรจึงจะเกิดการเลี้ยวเบนอย่างชัดเจน}
\begin{enumerate}
\item ก) ความยาวคลื่นมากกว่าช่องเปิด
\item ข) ความยาวคลื่นน้อยกว่าช่องเปิด
\item ค) ความถี่สูงมาก
\item ง) อัตราเร็วต่ำ
\end{enumerate}
\item \textbf{เสียงที่มนุษย์ได้ยินมีความถี่อยู่ในช่วงใด}
\begin{enumerate}
\item ก) 20 Hz - 20,000 Hz
\item ข) 20 kHz - 20,000 kHz
\item ค) 200 Hz - 2,000 Hz
\item ง) 2 Hz - 2,000 Hz
\end{enumerate}
\item \textbf{เมื่อคลื่นตกกระทบผิวสะท้อนแล้วทำมุม 40° กับเส้นแนวฉาก มุมสะท้อนมีค่าเท่าใด}
\begin{enumerate}
\item ก) 20°
\item ข) 40°
\item ค) 50°
\item ง) 80°
\end{enumerate}
\item \textbf{เมื่อคลื่นผิวน้ำเคลื่อนที่จากน้ำลึกไปน้ำตื้น ข้อใดเกิดขึ้น}
\begin{enumerate}
\item ก) ความถี่เพิ่มขึ้น
\item ข) ความยาวคลื่นลดลง
\item ค) อัตราเร็วลดลง
\item ง) ทิศเดิมไม่เปลี่ยน
\end{enumerate}
\item \textbf{คลื่นนิ่งเกิดจากคลื่นสองขบวนที่มีอะไรเหมือนกันแต่ต่างกัน}
\begin{enumerate}
\item ก) ความถี่เท่ากัน ทิศตรงข้ามกัน
\item ข) ความถี่ต่างกัน ทิศเดียวกัน
\item ค) ความยาวคลื่นเท่ากัน ทิศเดียวกัน
\item ง) ความถี่ต่างกัน ทิศตรงข้าม
\end{enumerate}
\item \textbf{เสียงที่มีความถี่ต่ำกว่า 20 Hz เรียกว่าอะไร}
\begin{enumerate}
\item ก) อัลตราซาวด์
\item ข) อินฟราซาวด์
\item ค) ไมโครซาวด์
\item ง) ไฮเปอร์ซาวด์
\end{enumerate}
\item \textbf{ในการทดลองช่องคู่ของยัง ระยะระหว่างปฏิบัพที่ 0 ถึงปฏิบัพที่ 3 วัดได้ 9 mm ระยะจากช่องถึงฉาย 1.5 m ความยาวคลื่นเป็นเท่าใด (d = 0.3 mm)}
\begin{enumerate}
\item ก) 450 nm
\item ข) 600 nm
\item ค) 750 nm
\item ง) 900 nm
\end{enumerate}
\item \textbf{ปรากฏการณ์ที่คลื่นสองขบวนเดินทางมาพบกันแล้วแอมพลิจูดรวมเป็นศูนย์ เรียกว่าอะไร}
\begin{enumerate}
\item ก) การแทรกสอดแบบเสริม
\item ข) การแทรกสอดแบบหักล้าง
\item ค) การเลี้ยวเบน
\item ง) การสะท้อน
\end{enumerate}
\item \textbf{คลื่นเสียงความถี่ 1000 Hz เดินทางในอากาศด้วยอัตราเร็ว 340 m/s เมื่อเข้าสู่น้ำ อัตราเร็วกลายเป็น 1500 m/s ความถี่และความยาวคลื่นในน้ำเป็นเท่าใด}
\begin{enumerate}
\item ก) f = 1000 Hz, λ = 1.5 m
\item ข) f = 1000 Hz, λ = 0.34 m
\item ค) f = 1500 Hz, λ = 1 m
\item ง) f = 441 Hz, λ = 3.4 m
\end{enumerate}
\item \textbf{คลื่นผิวน้ำในน้ำตื้น เมื่อเข้าหาชายฝั่งที่ลาดเอียง คลื่นจะเปลี่ยนแปลงอย่างไร}
\begin{enumerate}
\item ก) อัตราเร็วคงที่ ความยาวคลื่นคงที่
\item ข) อัตราเร็วลดลง ความยาวคลื่นลดลง ความถี่คงที่
\item ค) อัตราเร็วเพิ่มขึ้น ความยาวคลื่นเพิ่มขึ้น
\item ง) ความถี่เพิ่มขึ้น อัตราเร็วคงที่
\end{enumerate}
\item \textbf{หลักฮอยเกนส์ (Huygens' principle) กล่าวว่าอย่างไร}
\begin{enumerate}
\item ก) ทุกจุดบนขอบคลื่นเป็นแหล่งกำเนิดคลื่นใหม่
\item ข) คลื่นเดินทางเป็นเส้นตรงเสมอ
\item ค) แสงเดินทางเร็วกว่าเสียง
\item ง) คลื่นไม่สามารถเลี้ยวเบนได้
\end{enumerate}
\item \textbf{สเปกตรัมของคลื่นแม่เหล็กไฟฟ้าเรียงลำดับจากความถี่ต่ำไปสูง ข้อใดถูก}
\begin{enumerate}
\item ก) คลื่นวิทยุ → ไมโครเวฟ → อินฟราเรด → แสง → อัลตราไวโอเลต → รังสีเอกซ์ → รังสีแกมมา
\item ข) รังสีแกมมา → รังสีเอกซ์ → อัลตราไวโอเลต → แสง → อินฟราเรด → ไมโครเวฟ → คลื่นวิทยุ
\item ค) แสง → อินฟราเรด → ไมโครเวฟ → คลื่นวิทยุ
\item ง) คลื่นวิทยุ → อัลตราไวโอเลต → รังสีเอกซ์ → รังสีแกมมา
\end{enumerate}
\item \textbf{ในการทดลองแทรกสอดของแสงผ่านช่องคู่ ระยะระหว่างช่อง 0.4 mm ระยะจากช่องถึงฉาย 1.2 m ระยะระหว่างปฏิบัพที่ 2 ทั้งสองข้างของปฏิบัพกลางเป็น 6 mm ความยาวคลื่นเป็นเท่าใด}
\begin{enumerate}
\item ก) 500 nm
\item ข) 1000 nm
\item ค) 2000 nm
\item ง) 4000 nm
\end{enumerate}
\item \textbf{ทำไมเสียงจึงเดินทางได้เร็วกว่าในของแข็งมากกว่าในแก๊ส}
\begin{enumerate}
\item ก) เพราะของแข็งหนาแน่นกว่า
\item ข) เพราะโมเลกุลในของแข็งอยู่ชิดกันมากกว่า ส่งผ่านการสั่นได้เร็ว
\item ค) เพราะของแข็งมีอุณหภูมิสูงกว่า
\item ง) เพราะของแข็งมีความดันมากกว่า
\end{enumerate}
\item \textbf{คลื่นเสียงกำทินหนึ่งมีความเข้มเพิ่มขึ้น 4 เท่า ระดับเสียงจะเพิ่มขึ้นกี่เดซิเบล}
\begin{enumerate}
\item ก) 2 dB
\item ข) 4 dB
\item ค) 6 dB
\item ง) 8 dB
\end{enumerate}
\item \textbf{แสงสีขาวผ่านปริซึมแล้วแยกเป็นสีต่างๆ เรียกปรากฏการณ์นี้ว่าอะไร}
\begin{enumerate}
\item ก) การเลี้ยวเบน
\item ข) การแทรกสอด
\item ค) การกระจาย
\item ง) การสะท้อน
\end{enumerate}
\item \textbf{คลื่นแม่เหล็กไฟฟ้าทุกชนิดเคลื่อนที่ด้วยอัตราเร็วเท่ากันในสุญญากาศ อัตราเร็วนั้นเป็นเท่าใด}
\begin{enumerate}
\item ก) 3 × 10⁶ m/s
\item ข) 3 × 10⁸ m/s
\item ค) 3 × 10⁴ m/s
\item ง) 3 × 10⁷ m/s
\end{enumerate}
\item \textbf{คลื่นในข้อใดที่สามารถเดินทางผ่านสุญญากาศได้}
\begin{enumerate}
\item ก) คลื่นเสียง
\item ข) คลื่นในสปริง
\item ค) คลื่นแม่เหล็กไฟฟ้า
\item ง) คลื่นผิวน้ำ
\end{enumerate}
\item \textbf{แหล่งกำเนิดคลื่นสองแหล่งสั่นด้วยความถี่เท่ากัน 30 Hz อัตราเร็วคลื่น 3 m/s จุดบนแนวกลางระหว่างแหล่งที่เป็นบริเวณบัพหรือปฏิบัพ}
\begin{enumerate}
\item ก) เป็นบริเวณปฏิบัพเสมอ
\item ข) เป็นบริเวณบัพเสมอ
\item ค) เป็นได้ทั้งบัพและปฏิบัพขึ้นอยู่กับระยะ
\item ง) ไม่เกิดทั้งบัพและปฏิบัพ
\end{enumerate}
\item \textbf{ระดับเสียง 0 dB หมายความว่าอย่างไร}
\begin{enumerate}
\item ก) ไม่มีเสียง
\item ข) เสียงที่หูมนุษย์พึ่งได้ยิน
\item ค) เสียงดังมาก
\item ง) ความเข้มเสียงเป็น 1 W/m²
\end{enumerate}
\item \textbf{ถ้าเพิ่มความถี่ของคลื่นขึ้น 2 เท่า โดยอัตราเร็วคงที่ ความยาวคลื่นจะเป็นอย่างไร}
\begin{enumerate}
\item ก) เพิ่มขึ้น 2 เท่า
\item ข) ลดลงครึ่งหนึ่ง
\item ค) คงที่
\item ง) เพิ่มขึ้น 4 เท่า
\end{enumerate}
\item \textbf{ในการทดลองเลี้ยวเบนของคลื่นผ่านช่องเดี่ยว ความกว้างของช่องเท่ากับความยาวคลื่น ข้อใดจะเกิดขึ้น}
\begin{enumerate}
\item ก) ไม่มีการเลี้ยวเบน
\item ข) คลื่นแผ่ออกเป็นวงกลมเกือบทั้งหมด
\item ค) คลื่นสะท้อนกลับทั้งหมด
\item ง) เกิดเฉพาะแถบมืด
\end{enumerate}
\end{enumerate}
\end{document}